\section{Experimental Setup}
\label{sec:experimental_setup}

\subsection{Seven-system baseline comparison}
\label{sec:experimental_setup-model-configurations-and-controls}

The formal v5 comparison uses M0 direct answering, M2 MedRGAG with the local all-Llama
adaptation, M3 MA-RAG in its intrinsic configuration, M4 MedRAG--Llama, M5 MedRAG--Qwen,
M6 MedGENIE version 2, and M7 TC-RAG with the selected typed-control version 2. M0, M2,
M4, and M6 use Llama-3.1-8B-Instruct; M3, M5, and M7 use Qwen3-8B
\citep{dubey2024llama,yang2025qwen3}. Thus, seven configurations span six baseline
families and two underlying backbones. MA-RAG generates eight solver candidates per
round for at most four rounds, with early stopping on agreement. MedRGAG, MedGENIE,
and TC-RAG are the documented local adaptations, rather
than claims of exact reproduction of every published model and resource setting.

Each selected configuration covers the same 13,905 unique inputs and 15,616 scoring
mappings. Coverage combines new predictions with strictly matched historical reuse where
available: M0 reuses 8,877 and generates 5,028 predictions; M2 reuses 6,687 and generates
7,218; M3 reuses 3,690 and generates 10,215. M4--M7 each have 13,905 new predictions in
their selected complete run. Full coverage therefore means complete evaluated outputs,
not seven simultaneous runs from scratch. Invalid and unaccepted terminal outputs remain
in the common denominator.

The selected M4 baseline includes structured-output repair, whereas M5 retains its earlier
unconstrained output protocol. Corpora, upstream procedures, and budgets also vary across
families. Cross-system differences compare these concrete settings and cannot all be
assigned to an isolated algorithm or backbone effect. M1, an earlier MedRGAG initial
retrieval control, and the superseded i-MedRAG implementation are historical experiments,
not additional completed v5 configurations (Appendix~\ref{sec:appendix_development}).

\subsection{Seven-system terminal-head evaluation}
\label{sec:experimental_setup-terminal-heads}

We evaluate the frozen terminal heads on every member of R1, R2, and R3 for all seven
configurations. In R1, each configuration routes 1,780 diagnostic inputs to catalog
re-decision and 72 prescribing inputs to rules, retaining the other 1,014 answers. R3
contains 680 catalog inputs, 10 rule inputs, and 43 retained single-choice answers. R2
covers 55 of 79 inputs with rules and retains the remaining 24. Uncovered inputs remain
in their category denominators.

Diagnostic scores are generated with each configuration's actual original context and
backbone; earlier completed scores and exactly identical requests are reused as recorded.
The fact extractor sees the current question, options, and required fields, but no
method-specific retrieved history or initial answer. Consequently, its two backbone-specific
fact sets can be fixed and combined offline with seven original-answer sets. This is a
valid comparison of seven terminal systems, while the rule branch is not seven newly
executed upstream retrieval pipelines. No backbone is substituted to improve another
system's head result.

We compare the diagnostic head with single-pass catalog selection and uncorrected voting.
For R2, a constant ``none'' control intervenes on the same 55 inputs and preserves each
original answer on the other 24. Remaining-support correctness is reported alongside its
total. Additional same-facts execution, question-scope, and external-source evaluations
are kept within their original experimental scopes. The original formal baseline scores
and the development head scores remain separate tables.

\subsection{Two-configuration CF-MoA composition and verification}
\label{sec:experimental_setup-cf-moa}

The later six-arm experiment uses only the two MedRAG configurations M4 and M5, with
Qwen thinking disabled. Within each configuration, all six arms share the same frozen
full-v5 input scope and original contexts. It tests the composition of established heads
with further operations, followed by additional candidate verification; it does not
extend the six-arm experiment to all seven systems.

The ordinary F readout uses bfloat16 weights, temperature 0.7, top-p 1.0, a 2,048-token
output budget, and a 131,072-token context limit under the saved vLLM 0.8.5 profiles. The
F readout uses structured native output in both configurations. The existing catalog and
fact-extraction components retain their own shorter-context, temperature-zero profiles.
Exact per-component settings, prompts, code snapshots, and reuse bindings remain in the
released experiment package. The experiments use inference and deterministic tools; no new
model training is reported.

Table~\ref{tab:arms} summarizes the six evaluation arms.

\begin{table}[htbp]
\centering
\small
\caption{Six evaluation arms. C1 and B denote the same adopted operation core.}
\label{tab:arms}
\setlength{\tabcolsep}{5pt}
\begin{tabularx}{\linewidth}{@{}lXX@{}}
\toprule
Arm & Decision rule & Role in the comparison \\
\midrule
C0 & Original MedRAG configuration & Original-system reference \\
C1 / B & Adopted operation core & Strong system baseline \\
C2 & B plus candidate-only binary verification & Same-pool verification without saved rationales \\
C3 & B plus rationale-conditioned binary verification & Main research extension \\
C4 & B with old F votes plus three additional answers & Expanded-pool sampling control \\
C5 & B plus joint selection over the same candidates and rationales & Alternative same-pool selection control \\
\bottomrule
\end{tabularx}
\end{table}


C4 uses the predefined supplementary seeds 42, 43, and 44 on eligible ordinary-F
disagreements. A failed required supplementary answer leaves the control unavailable,
rather than silently converting it into a two-draw control. C5 scores a joint finite-code
choice under two opposing code assignments. C2, C3, and C5 preserve the existing pool.
Their comparisons share the original computation but have different additional request and
prompt costs; none is described as equal-cost merely because a common budget ceiling
exists.

Complete v5 is the scope of this later experiment. Earlier 121-input and 52-input
development panels informed its design and are not independent confirmation sets. The older 52-input
panel still lacks four original M5 C0 outputs; that historical gap does not affect the
complete full-v5 C0 results reported here. No new model calls or native grading are
performed for this manuscript; the figures and tables are rebuilt from frozen completed
results.

\subsection{Failure retention and computation
accounting}
\label{sec:experimental_setup-failure-retention-and-computation-accounting}

All arms retain the complete input and native-unit denominators. An unavailable answer
counts against performance under the frozen contract and is also reported explicitly. We
separate a valid but incorrect answer, an unavailable answer, and a verifier failure with
a retained F answer. Pool-trigger coverage includes the full scope and is not filtered by
subsequent success or gold correctness.

For the later composition experiment, per-method computation assigns the shared B computation once to each independently
considered B-based method, then adds that method's own incremental requests. This
estimates the attributable computation of a method assembled from the recorded components.
It is distinct from the physical ledger for the multi-arm experiment, where shared calls
were actually executed once. Exact cache reuse is not treated as free historical
computation. The upstream C0 and retrieval costs, and some historical event times, remain
unknown. Reported event seconds are not exclusive GPU hours, end-to-end latency, or energy
measurements.
